\documentclass[12pt]{article}
\usepackage[spanish,es-noshorthands]{babel}
\usepackage[utf8]{inputenc}
\usepackage[T1]{fontenc}
\usepackage{lmodern}
\usepackage{amsmath,amssymb,amsthm}
\usepackage{geometry}
\usepackage{enumitem}
\usepackage{multicol}
\usepackage{xcolor}
\usepackage{fancyhdr}
\geometry{margin=2.5cm}
\setlength{\parskip}{0.5em}
\setlength{\parindent}{0pt}
\pagestyle{fancy}
\fancyhf{}
\lhead{Solución -- Ayudantía 5}
\rhead{Trigonometría: Funciones y Aplicaciones}
\cfoot{\thepage}

\begin{document}
	
	\begin{center}
		{\LARGE \textbf{Solución completa -- Ayudantía 5}}\\[0.3em]
		{\large Trigonometría: Funciones y Aplicaciones}
	\end{center}
	
	\section*{Ejercicio 1: Transformaciones de sistemas de medición angular}
	Usamos las equivalencias:
	\[
	180^\circ = \pi \text{ rad}, \qquad 1^\circ = \frac{\pi}{180} \text{ rad}, \qquad 1\text{ rad} = \frac{180^\circ}{\pi}.
	\]
	
	\begin{enumerate}[label=\arabic*)]
		\item \(72^\circ\) a radianes:
		\[
		72^\circ\left(\frac{\pi}{180^\circ}\right)=\frac{72\pi}{180}=\frac{2\pi}{5}.
		\]
		
		\item \(-225^\circ\) a radianes:
		\[
		-225^\circ\left(\frac{\pi}{180^\circ}\right)=\frac{-225\pi}{180}=-\frac{5\pi}{4}.
		\]
		
		\item \(\dfrac{7\pi}{12}\) rad a grados:
		\[
		\frac{7\pi}{12}\left(\frac{180^\circ}{\pi}\right)=7\cdot 15^\circ=105^\circ.
		\]
		
		\item \(-\dfrac{11\pi}{5}\) rad a grados:
		\[
		-\frac{11\pi}{5}\left(\frac{180^\circ}{\pi}\right)=-11\cdot 36^\circ=-396^\circ.
		\]
	\end{enumerate}
	
	\textbf{Respuesta final:}
	\[
	72^\circ=\frac{2\pi}{5}, \qquad -225^\circ=-\frac{5\pi}{4}, \qquad \frac{7\pi}{12}=105^\circ, \qquad -\frac{11\pi}{5}=-396^\circ.
	\]
	
	\section*{Ejercicio 2: Ángulo entre las manecillas a las 4:00}
	A las 4:00, la manecilla minutera apunta al \(12\) y la horaria apunta al \(4\).
	
	Cada hora del reloj representa:
	\[
	\frac{360^\circ}{12}=30^\circ.
	\]
	
	Entonces, el ángulo entre ambas manecillas es:
	\[
	4\cdot 30^\circ=120^\circ.
	\]
	
	Como piden el ángulo menor, ese es justamente \(120^\circ\). Lo convertimos a radianes:
	\[
	120^\circ\left(\frac{\pi}{180^\circ}\right)=\frac{2\pi}{3}.
	\]
	
	\textbf{Respuesta final:}
	\[
	\boxed{\frac{2\pi}{3}\text{ rad}}
	\]
	
	\section*{Ejercicio 3: Razones trigonométricas exactas}
	
	\begin{enumerate}[label=\arabic*)]
		\item \(\sin(150^\circ)\)
		
		\(150^\circ\) está en el II cuadrante y su ángulo de referencia es \(30^\circ\). En el II cuadrante el seno es positivo.
		\[
		\sin(150^\circ)=\sin(30^\circ)=\frac{1}{2}.
		\]
		
		\item \(\cos(225^\circ)\)
		
		\(225^\circ=180^\circ+45^\circ\), está en el III cuadrante. En el III cuadrante el coseno es negativo.
		\[
		\cos(225^\circ)=-\cos(45^\circ)=-\frac{\sqrt{2}}{2}.
		\]
		
		\item \(\tan(300^\circ)\)
		
		\(300^\circ=360^\circ-60^\circ\), está en el IV cuadrante. En el IV cuadrante la tangente es negativa.
		\[
		\tan(300^\circ)=-\tan(60^\circ)=-\sqrt{3}.
		\]
		
		\item \(\sec\left(\dfrac{7\pi}{6}\right)\)
		
		\[
		\frac{7\pi}{6}=\pi+\frac{\pi}{6},
		\]
		está en el III cuadrante, donde el coseno es negativo.
		\[
		\cos\left(\frac{7\pi}{6}\right)=-\cos\left(\frac{\pi}{6}\right)=-\frac{\sqrt{3}}{2}.
		\]
		Entonces,
		\[
		\sec\left(\frac{7\pi}{6}\right)=\frac{1}{-\sqrt{3}/2}=-\frac{2}{\sqrt{3}}=-\frac{2\sqrt{3}}{3}.
		\]
		
		\item \(\csc\left(\dfrac{5\pi}{3}\right)\)
		
		\[
		\frac{5\pi}{3}=2\pi-\frac{\pi}{3},
		\]
		está en el IV cuadrante, donde el seno es negativo.
		\[
		\sin\left(\frac{5\pi}{3}\right)=-\sin\left(\frac{\pi}{3}\right)=-\frac{\sqrt{3}}{2}.
		\]
		Por tanto,
		\[
		\csc\left(\frac{5\pi}{3}\right)=\frac{1}{-\sqrt{3}/2}=-\frac{2}{\sqrt{3}}=-\frac{2\sqrt{3}}{3}.
		\]
		
		\item \(\cot\left(\dfrac{11\pi}{4}\right)\)
		
		Reducimos el ángulo módulo \(2\pi\):
		\[
		\frac{11\pi}{4}-2\pi=\frac{11\pi}{4}-\frac{8\pi}{4}=\frac{3\pi}{4}.
		\]
		Entonces,
		\[
		\cot\left(\frac{11\pi}{4}\right)=\cot\left(\frac{3\pi}{4}\right)=\frac{\cos(3\pi/4)}{\sin(3\pi/4)}=\frac{-\sqrt{2}/2}{\sqrt{2}/2}=-1.
		\]
		
		\item \(\sin\left(-\dfrac{5\pi}{6}\right)\)
		
		Usamos que el seno es función impar:
		\[
		\sin(-\theta)=-\sin(\theta).
		\]
		Entonces,
		\[
		\sin\left(-\frac{5\pi}{6}\right)=-\sin\left(\frac{5\pi}{6}\right)=-\frac{1}{2}.
		\]
	\end{enumerate}
	
	\textbf{Respuestas finales:}
	\[
	\sin(150^\circ)=\frac12,
	\quad
	\cos(225^\circ)=-\frac{\sqrt2}{2},
	\quad
	\tan(300^\circ)=-\sqrt3,
	\]
	\[
	\sec\left(\frac{7\pi}{6}\right)=-\frac{2\sqrt3}{3},
	\quad
	\csc\left(\frac{5\pi}{3}\right)=-\frac{2\sqrt3}{3},
	\quad
	\cot\left(\frac{11\pi}{4}\right)=-1,
	\quad
	\sin\left(-\frac{5\pi}{6}\right)=-\frac12.
	\]
	
	\section*{Ejercicio 4: Seis razones trigonométricas desde el punto \(P(-5,12)\)}
	Sea \(x=-5\), \(y=12\). Entonces:
	\[
	r=\sqrt{x^2+y^2}=\sqrt{(-5)^2+12^2}=\sqrt{25+144}=\sqrt{169}=13.
	\]
	
	Como \(x<0\) e \(y>0\), el ángulo está en el II cuadrante.
	
	Por definición:
	\[
	\sin\theta=\frac{y}{r}=\frac{12}{13},
	\qquad
	\cos\theta=\frac{x}{r}=-\frac{5}{13},
	\qquad
	\tan\theta=\frac{y}{x}=\frac{12}{-5}=-\frac{12}{5}.
	\]
	
	Y sus recíprocas:
	\[
	\csc\theta=\frac{13}{12},
	\qquad
	\sec\theta=-\frac{13}{5},
	\qquad
	\cot\theta=-\frac{5}{12}.
	\]
	
	\textbf{Respuesta final:}
	\[
	\boxed{\sin\theta=\frac{12}{13},\quad \cos\theta=-\frac{5}{13},\quad \tan\theta=-\frac{12}{5}}
	\]
	\[
	\boxed{\csc\theta=\frac{13}{12},\quad \sec\theta=-\frac{13}{5},\quad \cot\theta=-\frac{5}{12}}
	\]
	
	\section*{Ejercicio 5: Cálculo de \(E=\tan(\alpha)+\csc(\alpha)\)}
	Se sabe que:
	\[
	\cos(\alpha)=\frac13,
	\]
	con \(\alpha\) en el IV cuadrante. En el IV cuadrante, \(\cos\alpha>0\) y \(\sin\alpha<0\).
	
	Usamos la identidad pitagórica:
	\[
	\sin^2\alpha+\cos^2\alpha=1.
	\]
	Entonces,
	\[
	\sin^2\alpha=1-\left(\frac13\right)^2=1-\frac19=\frac89.
	\]
	Como \(\alpha\) está en el IV cuadrante:
	\[
	\sin\alpha=-\sqrt{\frac89}=-\frac{2\sqrt2}{3}.
	\]
	
	Ahora calculamos:
	\[
	\tan\alpha=\frac{\sin\alpha}{\cos\alpha}=\frac{-2\sqrt2/3}{1/3}=-2\sqrt2.
	\]
	
	Además,
	\[
	\csc\alpha=\frac{1}{\sin\alpha}=\frac{1}{-2\sqrt2/3}=-\frac{3}{2\sqrt2}=-\frac{3\sqrt2}{4}.
	\]
	
	Por tanto,
	\[
	E=\tan\alpha+\csc\alpha=-2\sqrt2-\frac{3\sqrt2}{4}.
	\]
	Llevando a común denominador:
	\[
	E=-\frac{8\sqrt2}{4}-\frac{3\sqrt2}{4}=-\frac{11\sqrt2}{4}.
	\]
	
	\textbf{Respuesta final:}
	\[
	\boxed{E=-\frac{11\sqrt2}{4}}
	\]
	
	\section*{Ejercicio 6: Deducción de \(\sec^2\alpha-\tan^2\alpha=1\)}
	Partimos de la identidad fundamental:
	\[
	\sin^2\alpha+\cos^2\alpha=1.
	\]
	
	Dividimos ambos lados por \(\cos^2\alpha\) (suponiendo \(\cos\alpha\neq 0\)):
	\[
	\frac{\sin^2\alpha}{\cos^2\alpha}+\frac{\cos^2\alpha}{\cos^2\alpha}=\frac{1}{\cos^2\alpha}.
	\]
	
	Simplificando:
	\[
	\tan^2\alpha+1=\sec^2\alpha.
	\]
	
	Restando \(\tan^2\alpha\) a ambos lados:
	\[
	1=\sec^2\alpha-\tan^2\alpha.
	\]
	
	Luego,
	\[
	\boxed{\sec^2\alpha-\tan^2\alpha=1.}
	\]
	
	\section*{Ejercicio 7: Justificación de \(\arcsin(1.5)\)}
	La función original es:
	\[
	f(x)=\sin(x).
	\]
	
	El recorrido del seno es:
	\[
	-1\leq \sin(x) \leq 1.
	\]
	
	Por lo tanto, la función inversa \(\arcsin(x)\) sólo está definida para valores de entrada que pertenezcan al intervalo:
	\[
	[-1,1].
	\]
	
	Como:
	\[
	1.5\notin[-1,1],
	\]
	se concluye que \(\arcsin(1.5)\) no está definida en los números reales. Por eso la calculadora arroja error matemático.
	
	\textbf{Conclusión:}
	\[
	\boxed{\arcsin(1.5)\text{ no existe en }\mathbb{R}.}
	\]
	
	\section*{Ejercicio 8: Funciones trigonométricas inversas}
	\subsection*{a) Evaluar \(\arcsin(\sin(5\pi/6))\)}
	Primero,
	\[
	\sin\left(\frac{5\pi}{6}\right)=\frac12.
	\]
	Entonces,
	\[
	\arcsin\left(\sin\left(\frac{5\pi}{6}\right)\right)=\arcsin\left(\frac12\right).
	\]
	
	Ahora bien, \(\arcsin(x)\) devuelve el ángulo principal en el intervalo:
	\[
	\left[-\frac\pi2,\frac\pi2\right].
	\]
	El ángulo de ese intervalo cuyo seno es \(1/2\) es:
	\[
	\frac\pi6.
	\]
	
	Por tanto,
	\[
	\arcsin\left(\sin\left(\frac{5\pi}{6}\right)\right)=\frac\pi6.
	\]
	
	La conjetura ``\(\arcsin(\sin\alpha)=\alpha\) siempre'' es \textbf{falsa} en general.
	
	\subsection*{b) Evaluar \(\cos(\arctan(-4/3))\)}
	Sea:
	\[
	\theta=\arctan\left(-\frac43\right).
	\]
	Entonces,
	\[
	\tan\theta=-\frac43.
	\]
	
	Como \(\arctan\) entrega ángulos principales en \((-\pi/2,\pi/2)\), \(\theta\) está en el IV cuadrante. Construimos un triángulo rectángulo con:
	\[
	\text{cateto opuesto}=-4, \qquad \text{cateto adyacente}=3.
	\]
	La hipotenusa es:
	\[
	r=\sqrt{(-4)^2+3^2}=\sqrt{16+9}=5.
	\]
	
	Así,
	\[
	\cos\theta=\frac{\text{adyacente}}{\text{hipotenusa}}=\frac35.
	\]
	Por lo tanto,
	\[
	\boxed{\cos\left(\arctan\left(-\frac43\right)\right)=\frac35.}
	\]
	
	\section*{Ejercicio 9: Distancia horizontal entre dos aeronaves}
	Sea \(x_A\) la distancia horizontal desde la base de la torre hasta la aeronave A, y \(x_B\) la distancia horizontal hasta la aeronave B.
	
	Como el ángulo de depresión es igual al ángulo de elevación correspondiente, se tiene:
	\[
	\tan(25^\circ)=\frac{50}{x_A}
	\qquad\Rightarrow\qquad
	x_A=\frac{50}{\tan(25^\circ)}=50\cot(25^\circ).
	\]
	
	Análogamente,
	\[
	\tan(15^\circ)=\frac{50}{x_B}
	\qquad\Rightarrow\qquad
	x_B=\frac{50}{\tan(15^\circ)}=50\cot(15^\circ).
	\]
	
	La separación horizontal entre ambas aeronaves es:
	\[
	D=x_B-x_A=\frac{50}{\tan(15^\circ)}-\frac{50}{\tan(25^\circ)}.
	\]
	Es decir,
	\[
	\boxed{D=50\left(\cot 15^\circ-\cot 25^\circ\right)}.
	\]
	
	Aproximando:
	\[
	D\approx 50(3.73205-2.14451)\approx 79.38\text{ m}.
	\]
	
	\textbf{Respuesta final:}
	\[
	\boxed{D=50\left(\cot 15^\circ-\cot 25^\circ\right)\text{ m}\approx 79.38\text{ m}}
	\]
	
	\section*{Ejercicio 10: Ángulo de elevación de una rampa}
	La rampa debe salvar un desnivel vertical de \(0.8\) m y su base horizontal máxima es \(6\) m. Entonces,
	\[
	\tan\theta=\frac{0.8}{6}=\frac{8}{60}=\frac{2}{15}.
	\]
	Por tanto,
	\[
	\theta=\arctan\left(\frac{2}{15}\right).
	\]
	
	Aproximando en grados:
	\[
	\theta\approx 7.59^\circ.
	\]
	
	Como la exigencia máxima legal es \(8^\circ\), se cumple que:
	\[
	7.59^\circ<8^\circ.
	\]
	
	\textbf{Conclusión:} la rampa \textbf{sí cumple} con la normativa.
	
	\textbf{Respuesta final:}
	\[
	\boxed{\theta=\arctan\left(\frac{2}{15}\right)\approx 7.59^\circ}
	\]
	
	\section*{Ejercicio 11: Altura de la montaña}
	Sea \(x\) la distancia horizontal desde el punto \(Q\) hasta la base de la montaña. Entonces, desde \(P\), la distancia horizontal hasta la base es \(x+800\).
	
	Sea \(h\) la altura de la montaña.
	
	Desde \(P\):
	\[
	\tan(30^\circ)=\frac{h}{x+800}.
	\]
	Como \(\tan 30^\circ=\frac{1}{\sqrt3}\), resulta:
	\[
	\frac{1}{\sqrt3}=\frac{h}{x+800}
	\qquad\Rightarrow\qquad
	h=\frac{x+800}{\sqrt3}.
	\]
	
	Desde \(Q\):
	\[
	\tan(45^\circ)=\frac{h}{x}.
	\]
	Como \(\tan45^\circ=1\), se tiene:
	\[
	1=\frac{h}{x}
	\qquad\Rightarrow\qquad
	h=x.
	\]
	
	Sustituyendo \(h=x\) en la primera ecuación:
	\[
	x=\frac{x+800}{\sqrt3}.
	\]
	Multiplicamos por \(\sqrt3\):
	\[
	x\sqrt3=x+800.
	\]
	Entonces,
	\[
	x(\sqrt3-1)=800.
	\]
	Despejando:
	\[
	x=\frac{800}{\sqrt3-1}.
	\]
	Racionalizamos:
	\[
	x=\frac{800}{\sqrt3-1}\cdot\frac{\sqrt3+1}{\sqrt3+1}
	=\frac{800(\sqrt3+1)}{3-1}
	=400(\sqrt3+1).
	\]
	Como \(h=x\), se concluye que:
	\[
	\boxed{h=400(\sqrt3+1)\text{ m}}.
	\]
	
	\section*{Resumen de respuestas}
	\begin{itemize}
		\item \textbf{Ej. 1:} \(\dfrac{2\pi}{5},\; -\dfrac{5\pi}{4},\; 105^\circ,\; -396^\circ\).
		\item \textbf{Ej. 2:} \(\dfrac{2\pi}{3}\) rad.
		\item \textbf{Ej. 3:} \(\dfrac12,\; -\dfrac{\sqrt2}{2},\; -\sqrt3,\; -\dfrac{2\sqrt3}{3},\; -\dfrac{2\sqrt3}{3},\; -1,\; -\dfrac12\).
		\item \textbf{Ej. 4:} \(\sin\theta=\dfrac{12}{13},\; \cos\theta=-\dfrac{5}{13},\; \tan\theta=-\dfrac{12}{5},\; \csc\theta=\dfrac{13}{12},\; \sec\theta=-\dfrac{13}{5},\; \cot\theta=-\dfrac{5}{12}\).
		\item \textbf{Ej. 5:} \(E=-\dfrac{11\sqrt2}{4}\).
		\item \textbf{Ej. 6:} \(\sec^2\alpha-\tan^2\alpha=1\).
		\item \textbf{Ej. 7:} \(\arcsin(1.5)\) no existe en \(\mathbb{R}\).
		\item \textbf{Ej. 8:} \(\arcsin(\sin(5\pi/6))=\dfrac\pi6\), \quad \(\cos(\arctan(-4/3))=\dfrac35\).
		\item \textbf{Ej. 9:} \(D=50(\cot15^\circ-\cot25^\circ)\approx 79.38\text{ m}\).
		\item \textbf{Ej. 10:} \(\theta=\arctan(2/15)\approx 7.59^\circ\), sí cumple.
		\item \textbf{Ej. 11:} \(h=400(\sqrt3+1)\text{ m}\).
	\end{itemize}
	
\end{document}